\section{Positive proportions of both horizontal directions}
\label{sec:signed-density}
The local velocity formula becomes a counting theorem by combining a
mixed phase average with a square mollifier. The phase average controls
where the horizontal component can be small; the mollifier controls
how strongly the derivative mass can concentrate there.

Fix a primitive character $\chi$ of conductor $q$, including $q=1$ for
zeta, and write $L_\chi(s)=L(s,\chi)$. All limits in this section keep
$\chi$ fixed. The translated family whose tangent is counted is
\[
 F_\chi(u,s)=\sum_{n\ge1}\chi(n)(n+u)^{-s}
 =q^{-s}\sum_{r=1}^q\chi(r)\zeta\left(s,\frac{r+u}{q}\right).
\]
The first identity holds for $\Re s>1$ and $|u|<1$; the finite
Hurwitz sum continues it near every nontrivial zero. At $u=0$ it
is $L_\chi(s)$ and $\partial_uF_\chi(0,s)=-sL_\chi(s+1)$,
so its simple-zero derivative is the velocity displayed below.
We use the zero-location theorem stated in the introduction,
including its assertion for the auxiliary Dirichlet characters in the
proof of Lemma~\ref{lem:dirichlet-additive}. Let $N_{\chi,T}$ count the
zero occurrences $\rho=1/2+i\gamma$, $T<\gamma\le2T$, with multiplicity.

\begin{theorem}\label{thm:positive-directions}
For either choice of sign, the number $n_{\chi,\pm}(T)$ of distinct
simple zeros in this interval satisfying
\[
 \pm\Re\frac{\rho L_\chi(\rho+1)}{L_\chi'(\rho)}>0
\]
obeys
\[
 \liminf_{T\to\infty}\frac{n_{\chi,\pm}(T)}{N_{\chi,T}}
       \ge10^{-14}.
\]
The numerical constant is independent of the fixed character. The height
from which a bound with any smaller constant holds may depend on it.
\end{theorem}

\subsection{A mixed phase average}
Let $\epsilon_\chi$ be the root number, $e_\chi\in\{0,1\}$ its parity,
and
\[
 Q_\chi(s)=(q/\pi)^{(s+e_\chi)/2}\Gamma((s+e_\chi)/2),\qquad
 \Delta_\chi(s)=\epsilon_\chi Q_\chi(1-s)/Q_\chi(s).
\]
Choose a continuous phase $\theta_\chi$ for which
$Z_\chi(t)=e^{i\theta_\chi(t)}L_\chi(1/2+it)$ is real. Then
$\Delta_\chi(1/2+it)^{-1}=e^{2i\theta_\chi(t)}$.
For $\sigma>1$, let $\phi_{\chi,\sigma}(t)$ be the imaginary part
of the Euler logarithm of $L_\chi(\sigma+it)$.

\begin{lemma}\label{lem:mixed-phase}
With $\delta_t=\arctan(1/(2t))$ and
$\psi_\chi(t)=\theta_\chi(t)+\phi_{\chi,3/2}(t)-\delta_t$, one has
\[
 \sum_{T<\gamma\le2T}\cos^2\psi_\chi(\gamma)
       =(\tfrac12+o_\chi(1))N_{\chi,T}.
\]
At a simple zero the absolute horizontal component of its velocity,
divided by the absolute velocity, is $|\cos\psi_\chi(\gamma)|$.
\end{lemma}

\begin{proof}
For every fixed $x>0$ we first prove
\begin{equation}\label{eq:twisted-hardy}
 \sum_{0<\gamma\le T}e^{2i\theta_\chi(\gamma)}x^{i\gamma}
 =-\frac{i^{e_\chi}\epsilon_\chi^{-1}}{\sqrt{qx}}
       \sum_{n\le qxT/(2\pi)}\chi(n)\Lambda(n)e^{-2\pi in/(qx)}
       +O_{\chi,x}(\sqrt T\log^2T)=O_{\chi,x}(T).
\end{equation}
Here and below zero sums count multiplicity. Integrate
$(L_\chi'/L_\chi)(s)\Delta_\chi(s)^{-1}x^{s-1/2}/(2\pi i)$
on a rectangle with vertical sides $-3/2$ and $c=1+1/\log T$,
fixed bottom ordinate above $2\pi$, and top within one unit of $T$
at distance $\gg_\chi1/\log T$ from all ordinates. The zero count in
any unit interval is $O_\chi(\log T)$, from the completed logarithmic
derivative at $2+it$; its partial fractions give
$L_\chi'/L_\chi=O_\chi(\log^2T)$ on the top. Stirling gives
\[
 \Delta_\chi(\sigma+it)^{-1}
 =\epsilon_\chi^{-1}(qt/(2\pi))^{\sigma-1/2}
 e^{i(t\log(qt/(2\pi))-t+\pi(2e_\chi-1)/4)}(1+O_\chi(t^{-1})).
\]
The gamma factor on the left edge is $O(t^{-2})$, so that edge
has a bounded integral. The top contributes $O(\sqrt T\log^2T)$.
Expanding the right edge gives the coefficient $-\chi(n)\Lambda(n)$ and the phase
$2\pi u\log(qxu/(en))$, $u=t/(2\pi)$.

For clarity, the stationary-phase error is summable here. Put
$r=n/(qx)$. Outside $[r/2,2r]$, the first derivative of
$u\log(u/(er))$ has absolute value at least $\log2$;
integration by parts gives $O(\sqrt T\log T)$ after summation.
The range $r\le1$ is treated by splitting at $u=2$: the bounded
lower interval and the first-derivative test above it contribute
$O_{\chi,x}(\sqrt T\log T)$. Thus the stationary formula below is
used only for $r>1$.
On $[A,B]=[1,T/(2\pi)]\cap[r/2,2r]$, the stationary-phase
formula of Arias de Reyna \cite[Lemmas~2.2 and~2.5, p.~5]{arias}
is
\[
 \mathbf1_{A\le r\le B}r^c e^{i\pi/4-2\pi ir}
 +O\left(r^{c-1/2}+
 \frac{r^{c+1/2}}{|A-r|+\sqrt r}+
 \frac{r^{c+1/2}}{|B-r|+\sqrt r}\right).
\]
The first and lower-end errors are $O(\sqrt T\log T)$ using
$\sum_{n\le X}\Lambda(n)n^{-1/2}\ll\sqrt X$. Near the upper
endpoint the error is bounded by
$C\sqrt T\log T\sum_{n\asymp T}(|n-qxT/(2\pi)|+\sqrt T)^{-1}
\ll\sqrt T\log^2T$.
The phase constants multiply to $i^{e_\chi}$, proving the first
identity in \eqref{eq:twisted-hardy}; Chebyshev's prime bound proves
the second. Moving the top back to $T$ changes at most $O(\log T)$
unit-modulus summands.

The absolutely convergent series
\[
 \frac{L_\chi(\sigma+it)}{\overline{L_\chi(\sigma+it)}}
 =\sum_{m,n\ge1}\frac{\chi(n)\mu(m)\overline{\chi(m)}}{(mn)^\sigma}
                e^{it\log(m/n)}
\]
has absolute coefficient sum at most $\zeta(\sigma)^2$.
Apply \eqref{eq:twisted-hardy} at each fixed $x=m/n$, divide by
$N_\chi(T)\sim T\log T/(2\pi)$, and use dominated convergence.
This proves that the mixed mode
$e^{2i\theta_\chi(\gamma)}L_\chi(\sigma+i\gamma)/
\overline{L_\chi(\sigma+i\gamma)}$ has mean zero. Multiplication
by $e^{-2i\delta_\gamma}$ changes the mean by $o(1)$.
Subtracting the formulas at $2T$ and $T$, and taking real parts,
gives the assertion about $\cos^2\psi_\chi$.
At a simple zero,
$L_\chi'(\rho)=-ie^{-i\theta_\chi(\gamma)}Z_\chi'(\gamma)$
with real nonzero $Z_\chi'$. Since $\arg(i\rho)=\pi-\delta_\gamma$,
the velocity formula gives the last assertion exactly.
\end{proof}

\begin{corollary}[The size of the normal component]
\label{cor:unsigned-transversality}
Let $N_{\chi,s,T}$ count simple zeros in $(T,2T]$, and at
each such zero put $r_\rho=|\Re v_\rho|/|v_\rho|$.
Then
\[
 \liminf_{T\to\infty}\frac1{N_{\chi,T}}
       \sum_{\rho\text{ simple},\ T<\gamma\le2T}r_\rho^2
       \ge\frac16.
\]
For $0<\eta<1/2$, the number $G_{\chi,\eta}(T)$ of these
zeros with $r_\rho\ge\eta$ obeys
\[
 \liminf\frac{G_{\chi,\eta}(T)}{N_{\chi,T}}
 \ge\frac{1-4\eta^2}{6(1-\eta^2)},\qquad
 \liminf\frac{G_{\chi,\eta}(T)}{N_{\chi,s,T}}
 \ge\frac{1-4\eta^2}{4(1-\eta^2)}.
\]
In particular, the corresponding bounds at $\eta=1/4$ are
$2/15$ and $1/5$. For every fixed $\varepsilon>0$, the zeros
counted at a positive threshold also satisfy
$|\Re v_\rho|\gg_{\chi,\eta,\varepsilon}T^{1-\varepsilon}$.
The same limiting bounds hold with the interval $(0,T]$.
\end{corollary}

\begin{proof}
The degree-one pair-correlation theorem
\cite[Theorems~1.1--1.2]{rudnick-sarnak} gives
$N_{\chi,s,T}\ge(2/3-o(1))N_{\chi,T}$.
Here is its exact counting consequence. Take nonnegative
Schwartz tests $g$ with $g(0)=1$, Fourier support in $(-1,1)$,
and sine-kernel integrals tending to $1/3$, constructed in
Section~\ref{sec:survival}. A center of multiplicity $m\ge2$
contributes $m(m-1)\ge m$ to the ordered off-diagonal sum.
Subtracting the global pair sums at $2T$ and $T$ bounds the
nonsimple occurrences by $(1/3+o(1))N_{\chi,T}$; all other
terms are nonnegative.

Abbreviate $N=N_{\chi,T}$, $N_s=N_{\chi,s,T}$ and $G=G_{\chi,\eta}(T)$.
Lemma~\ref{lem:mixed-phase} gives
\[
 \sum_{\rho\text{ simple}}r_\rho^2
       \ge N_s-N/2+o(N),\qquad
 \sum_{\rho\text{ simple}}r_\rho^2
       \le G+\eta^2(N_s-G).
\]
Thus $G\ge N_s-N/[2(1-\eta^2)]+o(N)$, which proves
both fractions and the mean-square assertion. The absolutely
convergent Euler product bounds $|L(\rho+1)|$ below by
$\zeta(3/2)^{-1}$. The character growth bound and Cauchy's
formula in Section~\ref{sec:dirichlet-moments} give
$|L'(\rho)|\ll_{\chi,\varepsilon}T^\varepsilon$;
the exact velocity formula proves the speed estimate.
Finally sum over finitely many dyadic intervals and then let
the number of intervals tend to infinity. The omitted initial
segment has arbitrarily small relative zero count, proving
the versions with cutoff $T$.
\end{proof}

\subsection{The square mollifier and its two constants}
Let $d_{-1/2}(n)$ be the coefficients of $\zeta(s)^{-1/2}$ in
$\Re s>1$. They are real, multiplicative and bounded in modulus by
one. Fix $0<\nu<1/4$ and a real polynomial $P$ satisfying $P(1)=1$
and having a zero of order at least sixteen at zero. Define
\[
 B_{\chi,T}(s)=\sum_{n\le T^\nu}d_{-1/2}(n)\chi(n)n^{-s}
                 P\left(\frac{\log(T^\nu/n)}{\nu\log T}\right).
\]
Write $M=T^\nu$ and
$b_T(n)=d_{-1/2}(n)P(\log(M/n)/\log M)\mathbf1_{n\le M}$.
For $k\ge1$, the binomial expansion gives
\[
 d_{-1/2}(p^k)=-\frac{\binom{2k}{k}}{4^k(2k-1)}.
\]
With $C_P=\max(1,\sup_{[0,1]}|P|)$, the exact squared coefficient
and its bound are therefore
\[
 a_T(n)=\sum_{uv=n}b_T(u)b_T(v),\qquad
 |a_T(n)|\le C_P^2\tau_2(n),\qquad a_T(n)=0\quad(n>M^2).
\]
Thus the square has length $T^{2\nu}<T^{1/2}$. The contour estimates
below apply this divisor envelope to the literal coefficients.
Write $L=\log T$.

\begin{proposition}\label{prop:square-moments}
There are real constants $a=a(P,\nu)$ and $b=b(P,\nu)$ such that
\begin{align}
 \sum_{T<\gamma\le2T}L_\chi'(\rho)B_{\chi,T}(\rho)^2
       &=(a+o_\chi(1))N_{\chi,T}L,\label{eq:square-first}\\
 \sum_{T<\gamma\le2T}|L_\chi'(\rho)B_{\chi,T}(\rho)^2|^2
       &=(b+o_\chi(1))N_{\chi,T}L^2.\label{eq:square-second}
\end{align}
They are independent of the fixed character. In addition, putting
\[
 K_{\rho,\chi}=\frac{\Re\{\rho L_\chi(\rho+1)
                  \overline{L_\chi'(\rho)}\}}{\gamma},
\]
one has
\begin{equation}\label{eq:square-cancel}
 \sum_{T<\gamma\le2T}K_{\rho,\chi}|B_{\chi,T}(\rho)|^2
       =o_\chi(N_{\chi,T}L).
\end{equation}
If $I_P=\int_0^1P(t)dt$ and $J_P=\int_0^1(1-t)P(t)dt$, then
\[
 a(P,\nu)=\tfrac12+\nu I_P+\tfrac{\nu^2}{4}(I_P^2-J_P).
\]
The constant $b$ is the finite integral defined below.
\end{proposition}

The discrete moment formulas used to prove this proposition are given,
with their contour proofs for all fixed primitive characters and the
displayed divisor envelopes, in Section~\ref{sec:dirichlet-moments}.
For zeta, the main terms agree with Benli--Elma--Ng
\cite[Theorem~1.1, equations~(9)--(14), pp.~3--4]{benli-elma-ng}
and Heap--Li--Zhao \cite[Theorem~2, equations~(5)--(9), pp.~3--5]{heap-li-zhao}.
The coefficient hypotheses of the first reference include its
submultiplicativity condition~(7); the contour derivation here proves
the needed formula directly from the divisor envelope above.
The second moment counts zero occurrences with multiplicity.
We give the kernel evaluation and
the rational inequality here, since they determine the positive counting
proportion.

Introduce variables $s_1,s_2,t_1,t_2$ and formal series
\[
 E_L(x)=\prod_{i=1}^2(1+x/s_i)^{1/2}=\sum l_kx^k,
 \qquad E_R(x)=\prod_{j=1}^2(1+x/t_j)^{1/2}=\sum r_kx^k.
\]
Define the following two finite polynomials:
\begin{align*}
 j={}&l_1r_1+l_1r_2+\tfrac12l_1+l_2r_1-l_3
                       +\tfrac12r_1-r_3+\tfrac13,\\
 i={}&l_1r_1^2+l_1r_1r_2+l_1r_3-\tfrac13l_1
       +l_2r_1^2-l_2r_1-l_2r_2-\tfrac12l_2-2l_3r_1+l_4\\
    &+\tfrac12r_1^2+r_1r_2-r_1r_3+\tfrac16r_1+r_2^2
                  +\tfrac12r_2-r_4-\tfrac18.
\end{align*}
Expand
\[
 j+2i=\sum_{\mathbf k}c_{\mathbf k}s_1^{-k_1}s_2^{-k_2}
                                      t_1^{-k_3}t_2^{-k_4}.
\]
This specifies sixty-two rational coefficients, of total degree at most
four. For $P(z)=\sum p_jz^j$ set
\[
 \Psi_k^P(z)=\sum_{j\ge1}p_j
        \frac{\Gamma(j+1)}{\Gamma(j+k+1/2)}z^{j+k-1/2}.
\]
On the domain
\[
 \mathcal D=\{x_{ij}\ge0:\ \sum_jx_{ij}\le1,\ \sum_ix_{ij}\le1\}
\]
put $z_1=1-x_{11}-x_{12}$, $z_2=1-x_{21}-x_{22}$,
$z_3=1-x_{11}-x_{21}$ and $z_4=1-x_{12}-x_{22}$. Then
\begin{equation}\label{eq:square-b}
 b(P,\nu)=\frac1{\Gamma(1/4)^4}
 \int_{\mathcal D}\prod_{i,j}x_{ij}^{-3/4}
   \sum_{\mathbf k}c_{\mathbf k}\nu^{|\mathbf k|-1}
            \prod_{l=1}^4\Psi_{k_l}^P(z_l)\,\prod_{i,j}dx_{ij}.
\end{equation}
All singularities of the weight are integrable.

\begin{proof}[Evaluation of the constants in Proposition~\ref{prop:square-moments}]
The first-moment residue calculation in Section~\ref{sec:dirichlet-moments}
gives, writing $a(n)$ for the coefficients of $B_{\chi,T}^2$, the leading
expression
\[
 \tfrac12L^2+\sum_p\frac{a(p)}p
       (-L\log p+\tfrac12\log^2p)
       +\sum_{p<r}\frac{a(pr)}{pr}\log p\log r+O_{\chi,P}(L).
\]
For $u=\log p/\log(T^\nu)$ and $v=\log r/\log(T^\nu)$, the
coefficients are
\[
 a(p)=-P(1-u)\mathbf1_{u\le1},\qquad
 a(pr)=\tfrac12\{P(1-u-v)\mathbf1_{u+v\le1}
       +P(1-u)P(1-v)\mathbf1_{u\le1,\ v\le1}\}.
\]
The measure $(\log Y)^{-1}\sum_{p\le Y}(\log p)/p\,
\delta_{\log p/\log Y}$ tends weakly to Lebesgue measure on $[0,1]$,
by the prime number theorem and partial summation. Diagonal prime
pairs and repeated prime powers contribute $o(L^2)$. The prime
integral is $\nu I_P-\nu^2J_P/2$; the pair integral is
$\nu^2(J_P+I_P^2)/4$. Their sum gives $a$.

We describe every limiting operation for the second constant. Use four
inverse Mellin variables for the four factors of $B$. Scaling each
variable by $L$ gives the weight
\[
 W(s)=e^{\nu s}\sum_j\frac{p_jj!}{\nu^j s^{j+1}},\qquad\Re s=1/\nu.
\]
The gcd sum in the continuous part of the discrete moment has Euler
product
\begin{equation}\label{eq:square-euler-F}
 C(u,v;\alpha,\beta)
 \frac{\prod_{i,j}Z_q(1+u_i+v_j)^{1/4}}
 {\prod_i Z_q(1+u_i+\beta)^{1/2}
       \prod_j Z_q(1+v_j+\alpha)^{1/2}},
\end{equation}
where $Z_q(s)=\zeta(s)\prod_{p\mid q}(1-p^{-s})$. The shifted
right-edge part has Euler product
\begin{align}\label{eq:square-euler-G}
 D(u,v;\alpha,\beta,\gamma)&
 \frac{Z_q(1+\alpha+\beta)Z_q(1+\beta+\gamma)}{Z_q(1+\beta)}
 \frac{\prod_{i,j}Z_q(1+u_i+v_j)^{1/4}}
       {\prod_iZ_q(1+u_i+\beta)^{1/2}}\nonumber\\
 &\hspace{10mm}\cdot\prod_j
 \frac{Z_q(1+v_j)^{1/2}}
 {Z_q(1+v_j+\alpha)^{1/2}Z_q(1+v_j+\gamma)^{1/2}}.
\end{align}
To verify these identities, pair one prime from the left coefficient
with one from the right. For the first product the linear terms are
$-\frac12\sum_i p^{-1-u_i-\beta}
-\frac12\sum_jp^{-1-v_j-\alpha}
+\frac14\sum_{i,j}p^{-1-u_i-v_j}$.
For the second, use the balanced relation $hm=kn$, with coefficient
$\mu*n^{-\alpha}*n^{-\gamma}$ on $m$ and $n^{-\beta}$ on $n$;
the same pairing gives precisely the displayed factors.
After removing these linear terms, the local errors are uniformly
$O(p^{-3/2})$ when the real parts are sufficiently near zero.
Thus $C,D$ converge absolutely there. At the origin the two half-power
series on each coefficient side combine to $\mu$. The first local
gcd sum is $1-1/p$; in the second, the combined left and right
coefficients are both $\delta_1$. These observations prove
$C(0)=D(0)=1$, including the prime powers. Primes dividing $q$ are
omitted from both sides.

Put $\alpha=A/L$, $\beta=B/L$, $\gamma=C_1/L$, on distinct fixed
small Cauchy circles in $A,B,C_1$. Every coefficient-variable zeta
argument remains to the right of one. On the complete Mellin lines,
$|Z_q(1+\epsilon+it)|+|Z_q(1+\epsilon+it)^{-1}|=O_q(L)$ for
$\epsilon\asymp1/L$. The two Euler expressions above are therefore
$O_q(L^3)$ and $O_q(L^6)$, respectively. The four weights decay at
least as $(1+|\Im s|)^{-17}$. Removing any region
$|\Im s|>\sqrt L$ costs $O_q(L^{-5})$ and $O_q(L^{-2})$, respectively;
these are smaller than the respective scales $L^{-1}$ and one.
On the remaining box use
$Z_q(1+z)=E_qz^{-1}(1+O_q(z))$, $E_q=\varphi(q)/q$.
The first expression has total exponent $-1$ and its explicit pole
in the continuous part cancels $E_q^{-1}$; the second expression has
total exponent zero. Thus no conductor factor remains.

Write
\[
 F_0(s,t)=\frac{(s_1s_2t_1t_2)^{1/2}}
                    {\prod_{i,j}(s_i+t_j)^{1/4}}.
\]
The two scaled limits are $F_0E_L(B)E_R(A)$ and
\[
 R(A,B,C_1)=F_0\frac{B}{(A+B)(B+C_1)}E_L(B)E_R(A)E_R(C_1).
\]
They have polynomial growth, while $s_i+t_j$ stays in a right
half-plane. The Mellin weights give an integrable majorant;
dominated convergence and Cauchy's formula justify every shift
derivative. The continuous part has normalized kernel
\[
 F_0\frac{E_L(B)E_R(A)-e^{-A-B}E_L(-A)E_R(-B)}{A+B},
\]
whose $AB$ coefficient is $F_0j$. The right-edge part is
$R(A,B,C_1)+e^{-A-B}R(-B,-A,C_1)+e^{-B-C_1}R(A,-C_1,-B)$.
Dividing by $F_0$, its numerator over
$(A+B)(B+C_1)(A-C_1)$ is
\begin{align*}
 &B(A-C_1)E_L(B)E_R(A)E_R(C_1)\\
 &\quad-A(B+C_1)e^{-A-B}E_L(-A)E_R(-B)E_R(C_1)\\
 &\quad+C_1(A+B)e^{-B-C_1}E_L(-C_1)E_R(A)E_R(-B).
\end{align*}
Each of the three denominator factors divides this numerator as a
convergent power series. Expansion through total degree six gives
$ABC_1$ coefficient $i$. Reflection gives the second copy, so the
full kernel is $F_0(j+2i)$.

Finally write
$(s_i+t_j)^{-1/4}=\Gamma(1/4)^{-1}\int_0^\infty
 x^{-3/4}e^{-x(s_i+t_j)}dx$. Absolute convergence justifies Fubini.
The single-variable inverse Laplace integral is
\[
 \frac1{2\pi i}\int W(s)s^{1/2-k}e^{-xs}ds
 =\nu^{k-1/2}\Psi_k^P(1-x/\nu)\mathbf1_{x\le\nu}.
\]
Scaling the four cross variables by $\nu$ gives a factor $\nu$ from
the measures and $\nu^{|\mathbf k|-2}$ from the four vertices.
This is exactly \eqref{eq:square-b}. The zero-location theorem
identifies the differentiated factors with a positive squared modulus,
which proves \eqref{eq:square-second}. The contour argument proving
\eqref{eq:square-cancel} appears in Section~\ref{sec:dirichlet-moments}.
\end{proof}

\subsection{A rational strict inequality}
We now prove a strict inequality for the constants, including the
smoothing required by the preceding Mellin argument. Put
$P_0(z)=2z-z^2$ and $\nu_0=1/4$. The finite functional gives
$a_0=1543/2304$ and
\[
 \Psi_k^{P_0}(z)=\frac{2z^{k+1/2}}{\Gamma(k+3/2)}
                         \left(1-\frac z{k+3/2}\right).
\]
After factoring $\sqrt{z_1z_2z_3z_4}/\pi^2$ out of the kernel
sum in \eqref{eq:square-b}, denote its remaining polynomial by
$Q(S,D,E)$, where
$S=z_1+z_2=z_3+z_4$, $D=(z_1-z_2)^2$, $E=(z_3-z_4)^2$.
This defines $Q$ exactly through the sixty-two rational coefficients.
Its constant and linear-$D$ coefficients at $E=0$ are
\begin{align*}
 q_0(S)={}&\frac8{245581875}\bigl(776S^8+75340S^7+287460S^6
 -5504100S^5-11151525S^4\\
 &\hspace{23mm}+151420500S^3+43659000S^2
 -1746360000S+2619540000\bigr),\\
 q_D(S)={}&-\frac8{245581875}\bigl(5207S^6-140S^5+948840S^4
 +1451100S^3\\
 &\hspace{23mm}-49075950S^2+30492000S+261954000\bigr).
\end{align*}
For $0\le S\le2$ and $0\le D,E\le\min(S,2-S)^2$, one has
\begin{equation}\label{eq:kernel-majorant}
 Q(S,D,E)\le U(S,D,E):=q_0(S)+q_D(S)D-(2-3S/4)E.
\end{equation}
Here $q_D<0$, and $U\ge0$ on the larger domain
$0\le D\le\min(S,2-S)^2$, $0\le E\le(2-S)^2$.
We give a finite rational certificate. Set
\[
 A=-[D^2]Q(S,D,0)-[D^3]Q(S,D,0)\,D,\qquad
 B=-(2-3S/4)-\frac{Q(S,D,E)-Q(S,D,0)}E.
\]
The identity $U-Q=D^2A+EB$ is exact. After the substitutions
$(S,D,E)=(u,u^2v,u^2w)$ and $(2-u,u^2v,u^2w)$, all tensor
Bernstein coefficients of $A$ and $B$ on the unit cube are positive.
Their minimum values are
\[
 \begin{array}{c|cc}
 &S=u&S=2-u\\\hline
 A&5060984/81860625&64264/3898125\\
 B&4133699/13097700&38706253/491163750
 \end{array}
\]
The coefficients of $q_D(2u)$ are negative. The minimum Bernstein
coefficients of
$q_0(S)+q_D(S)\min(S,2-S)^2-(2-3S/4)(2-S)^2$
on the two halves are $167144539/6682500$ and
$1492032128/245581875$. These checks use the exact conversion
\[
 [B_{\mathbf k}]p=\sum_{\mathbf j\le\mathbf k}a_{\mathbf j}
       \prod_l\frac{\binom{k_l}{j_l}}{\binom{n_l}{j_l}},
 \qquad p(x)=\sum a_{\mathbf j}x^{\mathbf j},
\]
where $n_l$ is the degree in variable $l$. Bernstein basis functions
are nonnegative and sum to one, so these rational inequalities prove
\eqref{eq:kernel-majorant} on its entire domain.

Set $r_i=x_{i1}+x_{i2}$ and $x_{11}=r_1v_1$, $x_{21}=r_2v_2$.
Then $\sqrt{z_3z_4}\le S/2$, and positivity of $U$ allows the
column constraints to be dropped after replacing $Q$ by $U$.
The $v_i$ have normalized beta law $\operatorname{Beta}(1/4,1/4)$,
with $\mathbb EE=2(r_1^2+r_2^2)/3$. The remaining normalized row
variables $R_i$ have beta law $\operatorname{Beta}(1/2,3/2)$.
Their moments are $(1/2)_k/(2)_k$. Direct polynomial integration yields
\begin{align*}
 b(P_0,1/4)&\le\frac1{4\pi}\mathbb E\left[
 \frac S2\left(q_0(S)+q_D(S)(R_1-R_2)^2
 -(2-3S/4)\frac23(R_1^2+R_2^2)\right)\right]\\
 &=\frac{1505515846753}{536481792000\pi},\qquad S=2-R_1-R_2.
\end{align*}
The program \texttt{certify\_square\_mollifier} constructs the rational
kernel, records every Bernstein coefficient, and repeats this integration.
It uses exact rational arithmetic throughout.

Machin's identity and its alternating arctangent series give
$\pi>314159/100000$. Thus, writing
$B_0=1505515846753/[536481792000(314159/100000)]$,
\[
 2a_0^2-B_0=\frac{5046955580047}{1348324666343424}>0.
\]
To meet the high-order endpoint condition take
$P_K(z)=(2z-z^2)(1-(1-z)^K)^{15}$.
It vanishes to order sixteen, satisfies $P_K(1)=1$, and
$0\le P_K'\le32$. The derivative converges in $L^4[0,1]$ to
$P_0'$, by dominated convergence. The identities
\[
 \Psi_0^P(z)=\frac1{\sqrt\pi}\int_0^z\frac{P'(t)}{\sqrt{z-t}}dt,
 \qquad
 \Psi_k^P(z)=\frac1{\Gamma(k-1/2)}\int_0^zP(t)(z-t)^{k-3/2}dt
 \quad(k\ge1)
\]
and H\"older's inequality prove uniform convergence of every vertex
needed in \eqref{eq:square-b}. In particular
$\|\Psi_0^{P_K}-\Psi_0^{P_0}\|_\infty
\le3^{3/4}\|P_K'-P_0'\|_4/\sqrt\pi$.
The kernel weight is integrable, so $a$ and $b$ converge as well.
They are continuous in $\nu>0$. We may therefore fix once and for
all a rational $\nu<1/4$ sufficiently close to $1/4$, and then a
sufficiently large integer $K$, for which
\[
 a\ge a_*:=1543/2304-1/2000,\qquad
 b\le b_*:=B_0+1/2000<1.
\]
Both parameters remain fixed as $T$ tends to infinity.

\begin{proof}[Proof of Theorem~\ref{thm:positive-directions}]
Choose the fixed mollifier just constructed. Let $\eta=1/50$ and
$p_\eta=1/[2(1-\eta^2)]$. Lemma~\ref{lem:mixed-phase} shows that
the occurrences with $|\cos\psi_\chi|<\eta$ number at most
$(p_\eta+o(1))N_{\chi,T}$: their contribution to the average
$\cos^2\psi_\chi$ is at most $\eta^2$, and every other contribution
is at most one. Put $A_\rho=L_\chi'(\rho)B_{\chi,T}(\rho)^2$.
The first moment and the triangle inequality give
$\sum|A_\rho|\ge(a_*+o(1))N_{\chi,T}L$.
Cauchy and the second moment bound the mass on the excluded set by
$(\sqrt{b_*p_\eta}+o(1))N_{\chi,T}L$. Thus its complement carries
at least $(d+o(1))N_{\chi,T}L$, where
$d=a_*-\sqrt{b_*p_\eta}$.

The Euler product bounds
$\zeta(3/2)^{-1}\le|L_\chi(3/2+it)|\le\zeta(3/2)$ and the exact
phase formula imply
\[
 |K_{\rho,\chi}|\,|B_{\chi,T}(\rho)|^2
  =\frac{|\rho|}{\gamma}|L_\chi(\rho+1)|\,|A_\rho|
                  |\cos\psi_\chi(\gamma)|.
\]
The total absolute weighted mass is therefore at least
$(\eta d/\zeta(3/2)+o(1))N_{\chi,T}L$.
Equation~\eqref{eq:square-cancel} divides it equally between the
two signs, up to $o(N_{\chi,T}L)$. On either sign set the second
moment of this weighted statistic is at most
$(\zeta(3/2)^2b_*+o(1))N_{\chi,T}L^2$. A further application of
Cauchy gives for each sign the lower proportion
\[
 \frac{\eta^2d^2}{4\zeta(3/2)^4b_*}.
\]
Every nonzero summand belongs to a simple center: at a multiple zero
$L_\chi'(\rho)=0$. At a simple center its sign equals that of the
velocity, because multiplication by $|L_\chi'(\rho)|^2/\gamma$ is
positive. The exact rational computation gives
$a_*^2-b_*p_\eta>1/2000$ and $0<a_*,b_*<1$, hence $d>1/4000$.
Finally $\zeta(3/2)<3$, by the integral test. The last lower
proportion exceeds
$1/12960000000000>10^{-14}$, as claimed.
\end{proof}
